% Metódy inžinierskej práce

\documentclass[10pt,twoside,slovak,a4paper]{article}

\usepackage[slovak]{babel}
%\usepackage[T1]{fontenc}
\usepackage[IL2]{fontenc} % lepšia sadzba písmena Ľ než v T1
\usepackage[utf8]{inputenc}
\usepackage{graphicx}
\usepackage{amsmath}
\usepackage{booktabs}
\usepackage{tikz}
\usetikzlibrary{positioning,fit,arrows.meta}
\usepackage{url} % príkaz \url na formátovanie URL
\usepackage{hyperref} % odkazy v texte budú aktívne (pri niektorých triedach dokumentov spôsobuje posun textu)

\usepackage{cite}
%\usepackage{times}

\pagestyle{headings}

\title{Black Box: včasné varovanie študentov pred študijným preťažením\thanks{Semestrálny projekt v predmete Metódy inžinierskej práce, ak. rok 2026/27, vedenie: \textcolor{red}{Meno Priezvisko}}} % meno a priezvisko vyučujúceho na cvičeniach

\author{Kyryl Tkachenko, Tymofii Trembovetskyi, Alex Vyšný\\[2pt]
	{\small Slovenská technická univerzita v Bratislave}\\
	{\small Fakulta informatiky a informačných technológií}\\
	{\small \texttt{xtkachenko@stuba.sk}\quad\texttt{xtrembovetskyi@stuba.sk}\quad\texttt{xvysny@stuba.sk}}
	}

\date{\small 30. október 2026}



\begin{document}

\maketitle

\begin{abstract}
Vysokoškolskí študenti dostávajú termíny zadaní, testov a~skúšok z~viacerých
zdrojov a~že sa v~niektorom týždni nahromadí priveľa povinností, často zistia
až vtedy, keď je neskoro. Dlhodobé preťaženie pritom patrí medzi hlavné príčiny
stresu a~vyhorenia študentov. V~článku prinášame prehľad poznatkov o~študijnej
záťaži a~jej vnímaní, o~nástrojoch na sledovanie pohody a~o~systémoch včasného
varovania vo vzdelávaní aj mimo neho. Ukazujeme, že existujúce riešenia buď
termíny iba zobrazujú, alebo sledujú pohodu bez väzby na konkrétne povinnosti,
a~že takmer všetky ukladajú citlivé údaje mimo zariadenia používateľa. Na tomto
základe navrhujeme lokálnu aplikáciu Black Box, ktorá z~odhadov časovej
náročnosti termínov počíta týždennú záťaž, 7 až 14 dní vopred varuje pred
preťaženým týždňom a~po náročnom období zobrazí časovú os termínov, záťaže
a~zaznamenaných pocitov. Opisujeme model výpočtu záťaže inšpirovaný
sledovaním tréningovej záťaže športovcov, architektúru typu local-first
a~plán overenia s~merateľnými ukazovateľmi.
\end{abstract}



\section{Úvod}

Študent v~priebehu semestra pracuje naraz na piatich až siedmich predmetoch.
Termíny odovzdania zadaní sú v~informačnom systéme školy, zmeny termínov
testov prichádzajú e-mailom a~dohody o~tímových stretnutiach vznikajú
v~skupinových chatoch. Každý zdroj osve vyzerá zvládnuteľne. Problém nastáva,
keď sa termíny z~rôznych predmetov zbehnú do toho istého týždňa: študent si to
uvedomí až pár dní vopred, keď už nemá priestor prácu rozložiť.

Že takéto obdobia nie sú náhodné, ale opakujú sa v~predvídateľnom rytme
semestra, ukázal projekt StudentLife~\cite{Wang:StudentLife}. Autori počas
desiatich týždňov sledovali 48 študentov pomocou senzorov smartfónu
a~krátkych dotazníkov a~zistili, že stres stúpa a~spánok aj fyzická aktivita
klesajú práve v~obdobiach so zvýšeným počtom termínov. Táto štúdia bola
prvotným podnetom pre našu tému. Ak sa záťaž dá v~dátach spätne rozpoznať,
dá sa ju aj vopred predpovedať, a~to z~údajov, ktoré študent pozná: z~termínov
a~z~odhadu, koľko práce si vyžiadajú.

Aplikáciu, ktorú navrhujeme, sme nazvali Black Box podľa letovej
zapisovacej jednotky. Tá turbulencii nezabráni, ale zaznamená, čo sa stalo,
aby sa z~toho dalo poučiť. Naša aplikácia ide o~krok ďalej: okrem záznamu sa
snaží na turbulenciu upozorniť skôr, než nastane.

Článok je usporiadaný takto. V~časti~\ref{s:problem} zhŕňame poznatky
o~strese, vyhorení a~vnímaní študijnej záťaže. Časť~\ref{s:pristupy}
prehľadne porovnáva existujúce prístupy: kalendáre, nástroje na sledovanie
pohody, systémy včasného varovania vo vzdelávaní a~analogické riešenie zo
športovej vedy. V~časti~\ref{s:navrh} opisujeme návrh aplikácie Black Box
vrátane modelu výpočtu záťaže a~architektúry. Časť~\ref{s:overenie} predstavuje
plán overenia a~merateľné ukazovatele. Otvorené otázky rozoberáme
v~časti~\ref{s:diskusia} a~záverečné poznámky prináša časť~\ref{s:zaver}.



\section{Preťaženie študentov a~jeho príčiny} \label{s:problem}

\subsection{Stres a~vyhorenie}

Vyhorenie je podľa Maslachovej a~Leitera dlhodobá reakcia na chronické
pracovné stresory, ktorá má tri zložky: vyčerpanie, cynizmus (odstup od
práce) a~pocit neúčinnosti~\cite{Maslach:Burnout}. Schaufeli a~kol. ukázali,
že tento koncept sa dá preniesť aj na štúdium: štúdium je pre študenta prácou
so zadanými úlohami a~termínmi a~vyhorenie študentov sa dá spoľahlivo merať
v~rôznych krajinách~\cite{Schaufeli:Burnout}. Pre náš projekt je dôležité
najmä to, že vyhorenie nevzniká naraz, ale hromadí sa z~opakovaných období
preťaženia.

Rozsah problému dokladajú aj veľké prieskumy. Medzinárodný projekt Svetovej
zdravotníckej organizácie zistil, že približne 35\,\% študentov prvého
ročníka malo počas života a~31\,\% za posledných 12 mesiacov príznaky
aspoň jednej z~bežných duševných porúch~\cite{Auerbach:WMH}. Prehľad Pascoeovej
a~kol. zhŕňa, že akademický stres, najmä tlak termínov a~skúšok, znižuje
študijný výkon, motiváciu a~zvyšuje riziko predčasného ukončenia
štúdia~\cite{Pascoe:Stress}.

\subsection{Vnímaná a~skutočná záťaž}

Kember upozorňuje, že pocit preťaženia nezodpovedá jednoducho počtu hodín,
ktoré študent strávi učením~\cite{Kember:Workload}. Vnímanie záťaže závisí
aj od toho, ako je organizovaná výučba a~hodnotenie, ako náročné je učivo
pochopiť a~aká je motivácia študenta. Z~toho pre nás vyplývajú dva
dôsledky. Po prvé, nestačí sledovať celkový počet hodín za semester, ale jeho
rozloženie v~čase. Po druhé, objektívny odhad záťaže treba dopĺňať
subjektívnym údajom o~tom, ako sa študent cíti.

\subsection{Prečo študenti preťaženie nepredvídajú}

Aj keby mal študent všetky termíny pred očami, predpoveď záťaže mu sťažujú
dva dobre zdokumentované javy. Prvým je \emph{plánovací klam}: ľudia
systematicky podceňujú, koľko času im úloha zaberie, a~to aj vtedy, keď
vedia, že podobné úlohy v~minulosti trvali dlhšie~\cite{Buehler:Planning}.
Druhým je \emph{prokrastinácia}: metaanalýza Steela ukazuje, že úsilie sa
s~blížiacim sa termínom prudko zvyšuje a~že odkladanie je medzi študentmi
mimoriadne rozšírené~\cite{Steel:Procrastination}. Výsledkom je, že práca na
viacerých termínoch sa v~skutočnosti sústredí do niekoľkých dní pred nimi
a~prekryv termínov sa prejaví silnejšie, než by naznačoval rovnomerný plán.

Oba javy priamo ovplyvňujú návrh nášho modelu. Odhad hodín od používateľa
nemôžeme brať doslova, ale musíme ho kalibrovať podľa toho, ako sa jeho
odhady v~minulosti zhodovali so skutočnosťou. Prácu na termíne nemáme
rozkladať rovnomerne, ale s~väčšou váhou na dni tesne pred ním.



\section{Existujúce prístupy} \label{s:pristupy}

\subsection{Kalendáre a~plánovače úloh}

Bežné kalendáre a~zoznamy úloh dokážu zhromaždiť termíny z~viacerých zdrojov,
najmä vďaka štandardu iCalendar~\cite{RFC5545}, ktorý podporuje aj väčšina
univerzitných informačných systémov. Termín je však v~nich iba bod v~čase.
Kalendár nevie, že za odovzdaním projektu je dvadsať hodín práce a~za
priebežným testom päť, a~preto nedokáže povedať, ktorý týždeň bude ťažký.
Taký úsudok zostáva na študentovi, ktorý ho podľa
časti~\ref{s:problem} robí systematicky zle.

\subsection{Sledovanie pohody a~osobná informatika}

Na opačnej strane sú nástroje, ktoré sledujú, ako sa človek cíti.
Metodicky vychádzajú z~\emph{ekologického momentového hodnotenia} (EMA), pri
ktorom respondent opakovane hodnotí svoj stav v~prirodzenom prostredí,
namiesto toho, aby si ho neskôr spätne vybavoval~\cite{Shiffman:EMA}. Takýto
záznam je presnejší ako dotazník na konci obdobia, pretože ho neskresľuje
pamäť. StudentLife~\cite{Wang:StudentLife} kombinoval krátke EMA dotazníky
s~pasívnymi senzormi a~ukázal, že zmeny pohody súvisia s~priebehom semestra.
Išlo však o~výskumný nástroj: údaje sa odosielali na server a~vyhodnocovali
ich výskumníci, nie samotní študenti.

Li, Dey a~Forlizzi opísali systémy \emph{osobnej informatiky} piatimi
etapami: príprava, zber, integrácia, reflexia a~akcia~\cite{Li:PersonalInformatics}.
Podľa ich zistení používatelia najčastejšie zlyhávajú pri integrácii údajov
z~rôznych zdrojov a~pri reflexii, teda pri hľadaní súvislostí. Väčšina
aplikácií na sledovanie nálady končí pri zbere: ukáže graf nálady, ale
nepovie, čím bol spôsobený, lebo o~termínoch nič nevie.

\subsection{Systémy včasného varovania vo vzdelávaní}

V~oblasti učiacej sa analytiky vznikli systémy, ktoré varujú pred
neúspechom v~predmete. Známym príkladom je Course Signals na Purdueovej
univerzite, ktorý zo známok a~aktivity v~systéme na riadenie výučby
priraďoval študentom semafor rizika~\cite{Arnold:CourseSignals}. Tieto
systémy sú inštitucionálne: údaje spracúva škola, varovanie sa týka jedného
predmetu a~cieľom je študijný výsledok, nie celková záťaž študenta naprieč
predmetmi.

\subsection{Analógia zo športu: sledovanie tréningovej záťaže}

Najbližšie k~nášmu problému má prístup, ktorý sa používa v~športovej vede.
Tréneri sledujú záťaž športovca a~porovnávajú záťaž posledného týždňa
s~priemerom predchádzajúcich štyroch týždňov. Gabbett ukázal, že riziko
zranenia nestúpa ani tak s~celkovým objemom tréningu, ako s~jeho prudkým
nárastom oproti tomu, na čo je športovec zvyknutý~\cite{Gabbett:ACWR}. Tento
\emph{pomer akútnej a~chronickej záťaže} (ACWR) je jednoduchý a~zrozumiteľný,
čo je pre aplikáciu určenú bežnému používateľovi výhoda.

Impellizzeri a~kol. však upozorňujú na jeho metodické slabiny: pomer je
nestabilný pri malej chronickej záťaži, napríklad na začiatku obdobia,
a~jeho prahové hodnoty sa nedajú preberať bez overenia v~konkrétnom
kontexte~\cite{Impellizzeri:ACWR}. Z~analógie preto preberáme myšlienku
porovnávať záťaž s~individuálnym zvykom, ale nie konkrétne prahy. Tie
musíme overiť na študentoch.

\subsection{Súkromie údajov o~pohode}

Údaje o~tom, ako sa človek cíti, môžu súvisieť s~jeho zdravím. Všeobecné
nariadenie o~ochrane údajov ich v~takom prípade radí medzi osobitné kategórie
osobných údajov s~prísnejšou ochranou~\cite{GDPR} a~v~Slovenskej republike
ho dopĺňa zákon č.~18/2018~Z.\,z.~\cite{Zakon18}. Prax aplikácií pre duševné
zdravie je pritom nepriaznivá. Huckvale a~kol. zistili, že 33 z~36
skúmaných aplikácií pre depresiu a~odvykanie od fajčenia odosielalo údaje
tretím stranám, najmä reklamným a~analytickým službám, a~len časť z~nich to
uviedla v~zásadách ochrany súkromia~\cite{Huckvale:Privacy}.

Alternatívou je prístup \emph{local-first}, ktorý Kleppmann a~kol.
charakterizujú siedmimi ideálmi; medzi nimi sú práca bez pripojenia,
dlhodobá dostupnosť údajov a~to, že používateľ má nad svojimi údajmi
plnú kontrolu~\cite{Kleppmann:LocalFirst}. Pri osobných údajoch o~pohode je
najbezpečnejším riešením, keď sa zo zariadenia vôbec nedostanú.

\subsection{Zhrnutie}

Tabuľka~\ref{t:porovnanie} porovnáva jednotlivé skupiny prístupov podľa
piatich vlastností, ktoré považujeme pre náš problém za podstatné. Žiadny
z~nich nespája termíny, predpoveď záťaže a~záznam pohody v~nástroji, ktorý by
mal študent pod vlastnou kontrolou.

\begin{table}[tbh]
\centering
\caption{Porovnanie existujúcich prístupov (\textbullet{} áno,
$\circ$ čiastočne, -- nie).}
\label{t:porovnanie}
\small
\begin{tabular}{@{}lccccc@{}}
\toprule
Prístup & Termíny & Predpoveď & Pohoda & Reflexia & Lokálne \\
\midrule
Kalendáre, plánovače              & \textbullet & --          & --          & --          & $\circ$ \\
Aplikácie na sledovanie nálady    & --          & --          & \textbullet & $\circ$     & $\circ$ \\
StudentLife~\cite{Wang:StudentLife} & $\circ$   & --          & \textbullet & $\circ$     & -- \\
Course Signals~\cite{Arnold:CourseSignals} & --  & $\circ$     & --          & --          & -- \\
Black Box (návrh)                 & \textbullet & \textbullet & \textbullet & \textbullet & \textbullet \\
\bottomrule
\end{tabular}
\end{table}



\section{Návrh riešenia Black Box} \label{s:navrh}

Black Box je desktopová aplikácia s~tromi hlavnými funkciami: zhromaždenie
termínov, predpoveď týždennej záťaže s~varovaním a~spätný pohľad na
náročné obdobie. Používateľ pridá predmety a~termíny ručne alebo importom
súboru vo formáte iCalendar~\cite{RFC5545}. Ku každému termínu zadá odhad,
koľko hodín práce mu zaberie, a~typ termínu (zadanie, test, skúška), podľa
ktorého aplikácia predvolí dĺžku prípravy.

\subsection{Model týždennej záťaže} \label{s:model}

Pre termín $i$ označme $d_i$ jeho dátum, $h_i$ odhad hodín a~$w_i$ počet dní
prípravy. Odhad hodín najprv vynásobíme kalibračným koeficientom $\alpha$
(pozri ďalej) a~rozložíme ho do dní pred termínom tak, aby váha rástla
s~blížiacim sa termínom, ako to zodpovedá poznatkom
o~prokrastinácii~\cite{Steel:Procrastination}. Záťaž dňa $t$ je
\begin{equation}
L(t) = \alpha \sum_i h_i \, k_i(d_i - t), \qquad
k_i(x) = \frac{2\,(w_i - x)}{w_i\,(w_i + 1)} \quad \text{pre } x = 0, \ldots, w_i - 1,
\label{e:denna}
\end{equation}
pričom mimo tohto rozsahu je $k_i(x) = 0$. Jadro $k_i$ je lineárne klesajúce
so vzdialenosťou od termínu a~jeho súčet je 1, takže celkový počet hodín sa
zachová. Záťaž týždňa $T$ je súčtom dennej záťaže
$Z(T) = \sum_{t \in T} L(t)$.

Varovanie pred týždňom $T$ aplikácia vydá, ak je splnená aspoň jedna z~dvoch
podmienok:
\begin{equation}
Z(T) > c \qquad \text{alebo} \qquad
\frac{Z(T)}{\max\!\left(\frac{1}{4}\sum_{j=1}^{4} Z(T-j),\; c/2\right)} > \rho .
\label{e:varovanie}
\end{equation}
Prvá podmienka porovnáva záťaž s~kapacitou $c$, teda počtom hodín týždenne,
ktoré má študent popri výučbe na samostatnú prácu k~dispozícii; zadáva ju
sám. Druhá je obdoba pomeru akútnej a~chronickej záťaže~\cite{Gabbett:ACWR}
a~zachytí prudký nárast oproti predchádzajúcim štyrom týždňom. Dolná hranica
$c/2$ v~menovateli rieši nestabilitu pomeru na začiatku semestra, na ktorú
upozorňujú Impellizzeri a~kol.~\cite{Impellizzeri:ACWR}. Počiatočné hodnoty
$c = 20$~h a~$\rho = 1{,}5$ sú pracovné a~ich nastavenie je jedným z~cieľov
overenia. Model sa vyhodnocuje denne pre týždne, ktoré začínajú o~7 až 14
dní neskôr, aby mal študent na úpravu plánu aspoň týždeň.

\paragraph{Kalibrácia odhadov.}
Po uplynutí termínu sa aplikácia voliteľne opýta, koľko hodín práca
skutočne trvala ($s_i$). Kalibračný koeficient je medián pomerov
$s_i / h_i$ za posledných najviac desať termínov; kým ich používateľ
nemá aspoň päť, platí $\alpha = 1$. Medián sme zvolili namiesto priemeru,
aby jeden extrémny termín nezmenil predpoveď na celý semester. Kalibrácia
priamo reaguje na plánovací klam~\cite{Buehler:Planning}: ak študent
systematicky podceňuje, aplikácia jeho odhady zväčší.

\subsection{Denný záznam a~spätný pohľad}

Raz denne môže používateľ jedným kliknutím zaznamenať, ako sa cíti, na
stupnici od 1 do 5. Krátky a~opakovaný záznam vychádza z~metódy
EMA~\cite{Shiffman:EMA}; zámerne obmedzujeme jeho rozsah, aby ho používateľ
po pár týždňoch neprestal vypĺňať. Raz týždenne sa aplikácia navyše opýta,
či bol uplynulý týždeň preťažený. Táto odpoveď slúži ako referenčný údaj pri
overovaní presnosti varovaní (časť~\ref{s:overenie}).

Spätný pohľad zobrazí na spoločnej časovej osi tri vrstvy: termíny,
vypočítanú záťaž $Z(T)$ a~zaznamenané pocity. Podľa modelu osobnej
informatiky~\cite{Li:PersonalInformatics} tým aplikácia preberá etapy
integrácie a~reflexie, ktoré pri samostatných nástrojoch zostávajú na
používateľovi. Študent vidí, ktoré termíny sa na náročnom týždni podieľali
najviac, a~môže si nabudúce prácu na nich rozložiť skôr.

\subsection{Architektúra}

Aplikácia vzniká v~nástroji Tauri~\cite{Tauri}, ktorý spája jadro
v~jazyku Rust s~webovým používateľským rozhraním a~v~porovnaní s~nástrojmi
založenými na celom prehliadači vytvára výrazne menšie inštalačné balíky.
Údaje sú uložené v~lokálnej databáze SQLite~\cite{SQLite}, ktorá je
jedným súborom v~zariadení používateľa. Rozdelenie zodpovedností ukazuje
obr.~\ref{f:architektura}.

\begin{figure}[tbh]
\centering
\begin{tikzpicture}[
	font=\small,
	box/.style={draw, rounded corners, align=center, minimum height=8mm, inner sep=4pt},
	arr/.style={-{Stealth}, thick}
]
\node[box] (ics) {Import\\iCalendar};
\node[box, below=4mm of ics] (man) {Ručné\\zadanie};
\node[box, below=4mm of man] (mood) {Denný záznam\\pocitov};

\node[box, right=12mm of man, minimum width=30mm] (core) {Jadro (Rust)\\[2pt]\footnotesize model záťaže,\\\footnotesize kalibrácia, varovania};
\node[box, right=10mm of core] (db) {SQLite};
\node[box, above=6mm of core, minimum width=30mm] (ui) {Rozhranie (web)\\[2pt]\footnotesize kalendár, časová os};

\draw[arr] (ics) -- (core);
\draw[arr] (man) -- (core);
\draw[arr] (mood) -- (core);
\draw[arr, <->] (core) -- (db);
\draw[arr, <->] (core) -- (ui);

\node[draw, dashed, rounded corners, inner sep=6pt, fit=(ics)(mood)(core)(db)(ui),
	label={[font=\footnotesize]below:zariadenie používateľa (bez pripojenia na server)}] {};
\end{tikzpicture}
\caption{Architektúra aplikácie Black Box. Všetky údaje zostávajú
v~zariadení používateľa.}
\label{f:architektura}
\end{figure}

Aplikácia nemá žiadnu serverovú časť a~na svoju činnosť nepotrebuje
internet. Nadväzuje tým na ideály local-first~\cite{Kleppmann:LocalFirst}
a~zohľadňuje zistenia o~úniku údajov z~aplikácií pre duševné
zdravie~\cite{Huckvale:Privacy}. Používateľ môže svoje údaje kedykoľvek
exportovať do otvoreného formátu alebo ich úplne vymazať, čím aplikácia
napĺňa právo na prenosnosť a~vymazanie údajov podľa
GDPR~\cite{GDPR}. Synchronizácia medzi viacerými zariadeniami nie je
v~prvej verzii súčasťou návrhu.



\section{Plán overenia} \label{s:overenie}

Aby sme vedeli posúdiť, či aplikácia skutočne pomáha, stanovili sme
merateľné ukazovatele pre každú z~troch funkcií.

\paragraph{Presnosť varovaní.}
Aplikáciu bude počas ôsmich týždňov semestra používať skupina najmenej
30 študentov. Za \emph{skutočne preťažený} považujeme týždeň, ktorý študent
v~týždennej otázke označí ako preťažený alebo ktorého priemerný denný záznam
je najviac 2. Varovania porovnáme s~týmto referenčným údajom a~vypočítame
presnosť (podiel oprávnených varovaní) a~úplnosť (podiel zachytených
preťažených týždňov). Cieľom je presnosť aspoň 70\,\% a~úplnosť aspoň
60\,\% pri predstihu najmenej 7 dní. Ako východiskový model použijeme
jednoduché počítanie termínov v~týždni; náš model by ho mal prekonať
v~ukazovateli F1.

\paragraph{Prínos kalibrácie.}
Pri termínoch, ku ktorým študenti zadajú skutočný počet hodín, porovnáme
strednú absolútnu chybu odhadu pred kalibráciou a~po nej. Očakávame, že
kalibrácia chybu zníži, čo by potvrdilo vplyv plánovacieho
klamu~\cite{Buehler:Planning} aj v~našom kontexte.

\paragraph{Použiteľnosť.}
Na konci overenia vyplnia účastníci dotazník SUS~\cite{Brooke:SUS}. Cieľom
je priemerné skóre aspoň 71, ktoré podľa Bangora a~kol. zodpovedá slovnému
hodnoteniu použiteľnosti ako dobrej~\cite{Bangor:SUS}. Doplnkovým ukazovateľom je podiel dní,
v~ktorých študent vyplnil denný záznam; ak klesne pod 50\,\%, denný záznam
je potrebné zjednodušiť.

\paragraph{Technické požiadavky.}
Aplikácia musí správne načítať aspoň 95\,\% udalostí z~exportu
univerzitného informačného systému vo formáte iCalendar a~počas celého
overenia nesmie odoslať žiadne údaje mimo zariadenia, čo overíme
sledovaním sieťovej komunikácie.



\section{Diskusia a~otvorené otázky} \label{s:diskusia}

Návrh má niekoľko obmedzení, ktoré chceme počas projektu preskúmať.
Prvým je závislosť od odhadov používateľa. Kalibrácia síce koriguje
systematickú chybu, ale nepomôže, ak študent odhady nezadáva vôbec.
Zvažujeme preto predvolené odhady podľa typu termínu, ktoré by sa
postupne prispôsobovali.

Druhým obmedzením je meranie pohody jedným číslom. Je to vedomý kompromis
medzi presnosťou a~ochotou záznam vypĺňať: validované dotazníky vyhorenia
ako MBI-SS~\cite{Schaufeli:Burnout} majú 15 položiek a~na denné
použitie sa nehodia. Zostáva otvorené, či jednopoložkový záznam dostatočne
zachytí zmeny, ktoré nás zaujímajú.

Tretím je prenositeľnosť prahov. Ako ukazuje kritika ACWR~\cite{Impellizzeri:ACWR},
prah, ktorý funguje v~jednej skupine, nemusí fungovať v~inej. Je možné, že
hodnoty $c$ a~$\rho$ sa budú musieť prispôsobovať jednotlivcovi, podobne
ako kalibračný koeficient $\alpha$.

Napokon, lokálne uloženie údajov vylučuje akékoľvek spoločné
vyhodnocovanie. Na účely overenia preto účastníci exportujú anonymizovaný
súhrn ukazovateľov sami a~s~ich výslovným súhlasom.



\section{Záver} \label{s:zaver}

V~článku sme zhrnuli poznatky o~študijnom preťažení a~o~tom, prečo ho
študenti sami nepredvídajú. Ukázali sme, že dostupné nástroje riešia vždy
len časť problému: kalendáre poznajú termíny, ale nie záťaž, aplikácie na
sledovanie pohody poznajú pocity, ale nie ich príčiny, a~systémy včasného
varovania vo vzdelávaní sú inštitucionálne a~zamerané na známky.

Naším príspevkom je návrh aplikácie Black Box, ktorá tieto časti spája
lokálne v~zariadení študenta, a~model výpočtu týždennej záťaže, ktorý
preberá myšlienku pomeru akútnej a~chronickej záťaže zo športovej vedy
a~dopĺňa ju o~kalibráciu odhadov. Stanovili sme merateľné ukazovatele, podľa
ktorých posúdime, či aplikácia svoj cieľ spĺňa.

Otvorené zostáva najmä to, či sa prahy varovania dajú nastaviť všeobecne
a~či jednoduchý denný záznam stačí na zachytenie preťaženia. Na tieto
otázky odpovie overenie so skupinou študentov, ktoré je ďalším krokom
projektu.



%\acknowledgement{Ak niekomu chcete poďakovať\ldots}


% týmto sa generuje zoznam literatúry z obsahu súboru literatura.bib podľa toho, na čo sa v článku odkazujete
\bibliography{literatura}
\bibliographystyle{plain} % prípadne alpha, abbrv alebo hociktorý iný
\end{document}
